\subsection{Structured completion for coordinate-indexed arrays}
\label{sec:structured}
Arbitrary site-label permutations need not preserve geometric sparsity. A
special case does admit direct construction: unsigned coordinate bits in
row-major order, including a valid prefix of a binary-capacity rectangle.
\begin{lemma}[Exact coordinate distance and Walsh support]
\label{lem:structured}
For $\ell$-bit unsigned integers $a,b$, let $a',b'$ be their lower $\ell-1$
bits and $A=(-1)^{a_{\ell-1}}, B=(-1)^{b_{\ell-1}}$. Their distance obeys
\begin{equation}
 |a-b|=\frac{1+AB}{2}|a'-b'|
 +2^{\ell-2}(1-AB)+\frac{B-A}{2}(a'-b'),
 \qquad D_0=0,
 \label{eq:structured}
\end{equation}
where $D_0$ is the distance between two zero-bit integers. For $\ell\ge1$,
the Walsh expansion has exactly
\begin{equation}
 T_\ell=5\,2^\ell-4\ell-4
 \label{eq:structured-support}
\end{equation}
nonzero coefficients, including the constant.
\end{lemma}
\begin{proof}
If the high bits agree, the first term of Eq.~\eqref{eq:structured} remains
and the other two vanish. If they differ, their ordering fixes the sign of
$a-b$ and gives the displayed low-bit correction. This proves the distance
identity. Expand
$a'-b'=\sum_{i<\ell-1}2^{i-1}(Z_{b_i}-Z_{a_i})$ and accumulate identical
masks. The base case is $T_1=2$. For $\ell\ge2$, each old mask appears both
alone and multiplied by $AB$; the $4(\ell-1)$ new high--low two-bit masks
are distinct from both sets. Only the constant and $AB$ coefficients combine
with the middle term of Eq.~\eqref{eq:structured}. The new constant is the
positive mean distance. For $N=2^{\ell-1}$, the previous constant is
$(N^2-1)/(3N)$, so the new $AB$ coefficient is $-(2N^2+1)/(6N)$ and is also
nonzero. Hence $T_\ell=2T_{\ell-1}+4(\ell-1)$, which gives
Eq.~\eqref{eq:structured-support} by induction.
\end{proof}
This constructs a sparse list of coefficients without a distance table or LP.
Construction uses $O(2^\ell)$ arithmetic operations and entries.
Exact coefficients (integers divided by powers of two) and masks require $O(\ell)$ bits per entry, adding
an $O(\ell)$ factor when accounting for bit operations and storage.

For coordinate widths $k_x+k_y=k$, combine the two coefficient lists on their separate
bit supports. The bound is $O(2^{k_x}+2^{k_y})$ terms and
$O(k(2^{k_x}+2^{k_y}))$ independent-parity CX per net, or $O(k\sqrt m)$ for
balanced binary rectangles. This is a subclass guarantee, not an arithmetic
oracle for arbitrary shuffled coordinates. The extension preserves all
valid prefix distances; unused grid points are virtual, not newly legal sites.
A change in the number of physical sites changes the optimization problem,
whereas extending the binary truth table does not.

\subsection{From coefficients to circuits}
Each nonconstant Walsh term requires a rotation conditioned on a parity.
Computing and uncomputing every parity independently can repeat many CX
gates. Our primary synthesis method first combines terms from different
nets that act on the same global qubit support. It groups the remaining
terms by their highest-index qubit, computes parity onto that qubit, and
visits the control masks in reflected Gray-code order. Between rotations,
only controls that change are updated; each target is restored at the end
of its group. We refer to this method as Gray synthesis.

Two additional ordering methods provide comparisons. Fixed mask order
processes nets and their integer-valued masks in increasing order. Beam
search instead retains four candidate term sequences and extends each with
up to ten terms chosen by support overlap, ranking sequences by CX count
after adjacent cancellation. These methods retain separate terms for each
net. Their comparison with Gray synthesis therefore includes both term
aggregation and ordering. Algorithm~\ref{alg:phase} specifies the ordering
and tie rules. All three methods retain nonzero coefficients after the
numerical threshold; intentional approximation is evaluated separately.

\subsection{Alternative extensions and approximation}
\label{sec:extensions}
We compare five extensions while holding coordinates and site labels fixed. Besides zero and weighted-$\ell_1$, a virtual-coordinate extension
appends unused integer lattice points in row-major order inside the bounding
rectangle, extending rows above it as necessary. Manhattan distance then
specifies every binary-code pair. A nearest-code extension maps each invalid code
to a valid code of least Hamming distance, breaking ties by smaller code, and
uses the mapped pair's distance. Mean fill assigns the mean of all $m^2$ valid
ordered-pair distances whenever either label is invalid. All five extensions
retain every valid pair, including equal-site pairs. Virtual geometry does
not imply sparse coefficients when the original labels are shuffled, and zero
padding is not proved to maximize support or gate count.

For L1 and virtual completion, the approximate control tests coefficients in
increasing absolute magnitude, breaking ties by mask. It removes a term only
when the maximum valid-pair residual remains at most $\eta D$, where $D$ is
the valid-site diameter and $\eta\in\{.001,.01,.05\}$. We report the achieved
residual and the resulting assignment-error bound, not just the allowed
tolerance. This greedy truncation is not an optimal sparse approximation.
It measures circuit size against bounded objective error without retraining
the ansatz.

Symbolic beam and Gray circuits retain a free $\gamma$. Numeric comparisons
fix $\gamma=.371$ and compare Gray, Qiskit's generic diagonal synthesis and
the ZX-calculus optimizer PyZX 0.10.3, using full reduction followed by
circuit extraction. Each extracted
circuit must pass PyZX's equality check against its parsed input; independent
small Qiskit unitary checks test the conversion and all five extension policies.
OpenQASM angle conversion is checked separately. Fixing angles can enable
additional simplifications, so numerical and symbolic results are reported
separately.

